% ============================================================
%  Приложения
% ============================================================
\chapter{Подробности по данным}
\label{app:data-details}

\section*{Источник и состав выборки}

В качестве источника использованы дневные цены закрытия (close) с
Yahoo Finance.

\todoblock{Привести полный список тикеров и временной диапазон выборки.
Например: AAPL, MSFT, GOOG, AMZN, META; 2010-01-01 --- 2025-12-31.}

\section*{Скользящее окно реализованной корреляции}

Реализованная корреляция оценивается по матрице лог-доходностей в окне
$W = 21$ рабочих день:
\begin{equation}
  \hat\rho_{ij}(t)
  \;=\;
  \frac{\sum_{s=t-W+1}^{t}(r^{(i)}_s-\bar r^{(i)})(r^{(j)}_s-\bar r^{(j)})}
       {\sqrt{\sum (r^{(i)}_s-\bar r^{(i)})^{2}}\sqrt{\sum (r^{(j)}_s-\bar r^{(j)})^{2}}}.
\end{equation}

\section*{Оценка параметров модели Якоби}

Для модели Якоби~\eqref{eq:jacobi} параметры $(\kappa,\bar\rho,\sigma)$
оцениваются методом квази-максимального правдоподобия с использованием
дискретной аппроксимации Эйлера. Окно оценки~--- $252$ рабочих дня.

\chapter{Технические детали реализации}
\label{app:code}

Численные эксперименты выполнены на Python~3.11; для нейронных сетей
использован PyTorch, для финансовых моделей~--- библиотеки NumPy,
SciPy, statsmodels, а также собственная реализация формул из
\cref{ch:corr}.

\todoblock{Привести ссылку на репозиторий с кодом (если будет), таблицу
гиперпараметров обучения и среду исполнения.}
